\section{From columns to a four-row moment recursion}\label{sec:recursion}
Our objective is to express a moment on $N$ cards in terms of moments on $N/4$ cards. We first fix the operator conventions, then isolate the cost of the interaction between the four rows.

Positions are binary vectors. A permutation sends input positions to output positions, and $gh$ applies $h$ first. In each coordinate direction, the average of all optional edge swaps is an orthogonal projection $\Pi_i$. With chronological order $1,\ldots,d$,
\[
 T_N=\Pi_d\cdots\Pi_1,\qquad B_N=T_NT_N^*.
\]
To check physical time, let $R(x_1,\ldots,x_d)=(x_2,\ldots,x_d,x_1)$ and let $S_t$ be the switches in the first coordinate. A physical step is $RS_t$. If $Y_t$ is the accumulated permutation, then $X_t=R^{-t}Y_t$ satisfies
\[
 X_{t+1}=(R^{-t}S_{t+1}R^t)X_t.
\]
The conjugated directions visit all coordinates in a fixed cyclic order, and $R^d=I$. Relabeling that order as $1,\ldots,d$ gives exactly $T_N$ after $d$ steps.

These are group-algebra elements as well as operators in any representation. The Fourier transform convention is $\widehat\mu(\lambda)=\sum_g\mu(g)\rho_\lambda(g)$; convolution multiplies these matrices in the same order. For a matrix $A$, the unnormalized Schatten norm is $\|A\|_p^p=\Tr|A|^p$.

\begin{lemma}[Annihilation and the finite-size gap]\label{lem:gap}
Every sweep kills the sign representation and every shape with more than $N/2$ rows. For fixed dyadic $N\ge2$, its operator norm on each nontrivial irreducible representation is strictly less than one.
\end{lemma}
\begin{proof}
Each layer averages a subgroup $H\cong(S_2)^{N/2}$. The sign character averages to zero. Young's rule says that an irreducible with $H$-fixed vectors must dominate the partition $(2^{N/2})$, so it has at most $N/2$ rows. For the norm assertion, equality in a product of orthogonal projections on a unit vector would force that vector to be fixed by every layer. It would then be fixed by every coordinate-edge transposition. Edge transpositions of a connected graph generate its full symmetric group, whereas a nontrivial irreducible has no invariant vector.
\end{proof}

Take $N=4m$, using the last two coordinates as the row index. The first $d-2$ directions are independent sweeps within the four rows; the final two directions form a two-axis sweep independently in each of the $m$ columns. Let $M$ be the latter operator. Then
\begin{equation}\label{eq:split}
 B_N=MAM^*,\qquad A=B_m^{\otimes4}.
\end{equation}
The row group is $H=(S_m)^4$. The column group is $(S_4)^m$; its factors act on disjoint columns. Figure~\ref{fig:fourrows} records the two operations.

\begin{figure}[ht]
\centering
\begin{tikzpicture}[x=0.60cm,y=0.55cm,>=Stealth]
\foreach \i in {0,...,3}{\foreach \j in {0,...,7}{\fill (\j,-\i) circle (1.7pt);}}
\foreach \i in {0,...,3}{\draw[->,blue,thick] (-.6,-\i+.16)--(7.6,-\i+.16);}
\foreach \j in {0,...,7}{\draw[->,red,thick] (\j+.15,.6)--(\j+.15,-3.6);}
\node[blue,anchor=west] at (8,-1) {four sweeps on $m$ positions};
\node[red,anchor=west] at (8,-2) {$m$ sweeps on four positions};
\end{tikzpicture}
\caption{The four rows are indexed by the last two processed bits. The horizontal sweeps occur first; the vertical two-axis sweeps occur last. The number of columns shown is schematic.}\label{fig:fourrows}
\end{figure}

\paragraph{A positive trace for each row type.}
Fix an even integer $r\ge2$, and put
\[
 F_r(\lambda)=D_\lambda\Tr B_{|\lambda|}(\lambda)^r.
\]
We suppress the subscript $r$ while deriving the recursion. Set $X=A^{r/2}$ and $Y=(M^*M)^{r/2}$. The Araki--Lieb--Thirring inequality for positive operators, with exponent $r/2\ge1$ and outer power two, gives
\begin{equation}\label{eq:alt}
 F(\lambda)\le D_\lambda\Tr(XYXY)_\lambda
 =D_\lambda\|Y^{1/2}X_\lambda Y^{1/2}\|_{\HS}^2.
\end{equation}
Indeed $MAM^*$ and $A^{1/2}M^*MA^{1/2}$ have the same nonzero eigenvalues, after which the stated positive-matrix inequality applies \cite{araki1990}; see also the formulation in \cite[Theorem~1]{audenaert2007}.

The row operator $X=A^{r/2}$ is a positive element of the group algebra of
$H$. Let $\mathcal I(\lambda)$ be the set of tuples
$\boldsymbol\alpha=(\alpha_1,\ldots,\alpha_4)$, $\alpha_i\vdash m$,
occurring in $\Res_H V_\lambda$. For each tuple let
$X_{\boldsymbol\alpha}$ be $X$ multiplied by its central isotype
projection in the group algebra of $H$. This projection commutes with $X$,
so the summand is positive in every unitary representation. On $V_\lambda$
it acts on all copies of that row type, not on a chosen copy.
The Hilbert--Schmidt triangle inequality in \eqref{eq:alt} gives
\begin{equation}\label{eq:triangle}
 F(\lambda)^{1/2}\le
 \sum_{\boldsymbol\alpha\in\mathcal I(\lambda)}
 \bigl(D_\lambda\Tr(X_{\boldsymbol\alpha}Y
                         X_{\boldsymbol\alpha}Y)_\lambda\bigr)^{1/2}.
\end{equation}
We now estimate one term of this sum. In every irreducible representation,
$\Tr(X_{\boldsymbol\alpha}Y X_{\boldsymbol\alpha}Y)$ is the squared
Hilbert--Schmidt norm of $Y^{1/2}X_{\boldsymbol\alpha}Y^{1/2}$ and is
nonnegative. The regular representation of $S_N$ contains $D_\lambda$
copies of $V_\lambda$. Its trace therefore bounds the quantity inside
the corresponding square root in \eqref{eq:triangle}. Working in this
larger representation will allow us to express the bound as a probability.

\paragraph{From coefficients to an overlap probability.}
Fix the tuple $\boldsymbol\alpha$. Write
$X_{\boldsymbol\alpha}=\sum_{p\in H}\xi(p)p$, where
$\xi(p)=\prod_i\xi_i(p_i)$. Each row factor has Fourier block
$B_m(\alpha_i)^{r/2}$ on $\alpha_i$ and zero blocks elsewhere.
Plancherel on $S_m$ gives
\[
 m!\sum_{p_i}|\xi_i(p_i)|^2
 =D_{\alpha_i}\Tr B_m(\alpha_i)^r=F(\alpha_i).
\]
If one of these quantities is zero, the tuple contributes nothing.
Otherwise the squared coefficients define probability laws
$u_i(p_i)=|\xi_i(p_i)|^2/\sum|\xi_i|^2$ on the four row groups.
Write $u=\bigotimes_i u_i$ for their independent product.

Because $r/2$ is an integer, $Y=(M^*M)^{r/2}$ is convolution by a
probability law $\omega$ on the column group. Expanding the regular trace gives
\[
 N!\sum_{p,q\in H}\sum_{b,c}\xi(p)\xi(q)\omega(b)\omega(c)
                       \one_{\{p b q c=e\}}.
\]
Bound each coefficient product by
$|\xi(p)\xi(q)|\le(|\xi(p)|^2+|\xi(q)|^2)/2$.
The two square terms have equal sums: cyclically interchange $(p,b)$
and $(q,c)$ in the identity $pbqc=e$. For fixed $p,b,c$, there is
at most one possible $q$, and it lies in $H$ exactly when
$b^{-1}p^{-1}c^{-1}\in H$.
Self-adjointness gives $\xi(p^{-1})=\overline{\xi(p)}$, so $u$ is
invariant under inversion; the real probability coefficients of the
self-adjoint $Y$ have the same symmetry. We may thus invert each of
the three independent variables. Using the preceding Plancherel
normalizations yields
\begin{equation}\label{eq:coefficient-trace}
 \Tr_{\rm reg}(X_{\boldsymbol\alpha}Y
                         X_{\boldsymbol\alpha}Y)
 \le\prod_i F(\alpha_i)\,
 \frac{N!}{(m!)^4}\Pp\{bpc\in H\},
\end{equation}
where $p\sim u$, and $b,c$ are independent column permutations of law
$\omega$. The remaining problem is to bound this overlap probability
uniformly in the squared-coefficient laws $u_i$.

\paragraph{Two independent colorings.}
Let $a(z)$ be the row index of position $z$, and define
\[
 \eta(z)=a(c^{-1}z),\qquad \zeta(z)=a(bz).
\]
The condition $bpc\in H$ is equivalent to
$\zeta(pz)=\eta(z)$ for every position $z$: substitute $z=cw$ in
$a(bpcw)=a(w)$. In each column, either coloring uses the four colors
once. Different columns are independent, and the two colorings are
independent because $b$ and $c$ are.

For row $i$, define a transition matrix on all $4^m$ color strings by
\[
 U_i(x,y)=\Pp_{p_i\sim u_i}\{y(p_i z)=x(z)\text{ for every position }z
                         \text{ in row }i\}.
\]
Thus $U_i$ averages the permutation action of the row group on strings.
Conditional on the two colorings, the four row permutations are
independent. With $\eta_i,\zeta_i$ denoting their restrictions to row $i$,
\[
 \Pp\{bpc\in H\}=\E\prod_i U_i(\eta_i,\zeta_i).
\]
The matrices $U_i$ are stochastic and symmetric, since each $u_i$ is
invariant under inversion. The particular origin of $u_i$ gives the
additional positivity needed in \eqref{eq:interpolation}. Fourier
inversion writes
$\xi_i(g)=(D_{\alpha_i}/m!)\Tr(Q\rho_{\alpha_i}(g^{-1}))$ with
$Q=B_m(\alpha_i)^{r/2}\ge0$. Hence $\xi_i$, and then
$|\xi_i|^2$, are positive-definite functions. Their Fourier transforms
are positive semidefinite, so the average of $u_i$ in the color-string
representation is positive semidefinite as well.

The four-site law of $M$ sends every site exactly uniformly to the four sites. Its adjoint does too, and the same remains true of every positive power of $M^*M$. Moreover the support of $M^*M$ contains the identity and the edge transpositions of the square, which generate $S_4$. Its convolution powers therefore converge to uniform. Choose the even integer $r$ sufficiently large that the law of each column of $Y$ satisfies Theorem~\ref{thm:permanent}. Every larger even $r$ does too, once a tail threshold for convergence has been chosen.

\paragraph{Closing the recursion.}
The variables defining $(\eta_i,\zeta_i)$ consist of $2m$ independent
four-color column assignments. Apply Lemma~\ref{lem:tensorization} to
these assignments with the four nonnegative functions $U_i$, and then
\eqref{eq:interpolation} with $s=4^m$. This gives
\[
 \Pp\{bpc\in H\}
 \le\prod_i\left(s^{-2}\sum_{x,y}U_i(x,y)^{p_0}\right)^{1/p_0}
 \le s^{-4}\prod_i(\Tr U_i)^\theta,
 \quad\theta=1-1/p_0<\frac12.
\]
Since $N!/(m!)^4\le4^N=s^4$, the exponential factors cancel exactly. The only remaining inflation is a power of the four color-string traces. Write $c^\alpha_{\boldsymbol\beta}$ for the multiplicity of $\bigotimes_{i=1}^4V_{\beta_i}$ in the restriction of $V_\alpha$ to $\prod_{i=1}^4S_{|\beta_i|}$. Section~\ref{sec:multiplicities} proves
\begin{equation}\label{eq:trace-C}
 \Tr U_i\le C(\alpha_i),\qquad
 C(\alpha)=\sum_{\boldsymbol\beta}(c^{\alpha}_{\boldsymbol\beta})^2.
\end{equation}
The sum includes every weak four-part composition of $m$ and every
irreducible tuple for its Young subgroup. Combining this trace bound
with \eqref{eq:coefficient-trace} and \eqref{eq:triangle} gives
\begin{equation}\label{eq:recursion}
 F(\lambda)^{1/2}\le
 \sum_{\boldsymbol\alpha\in\mathcal I(\lambda)}
       \prod_{i=1}^4\left(F(\alpha_i)C(\alpha_i)^\theta\right)^{1/2}.
\end{equation}
All restriction multiplicities have been retained; they are accounted for by the regular trace and by $C(\alpha)$, rather than discarded during a choice of one copy.
